\documentclass[11pt]{article}
\usepackage[margin=1in]{geometry}
\usepackage{amsmath,amssymb,amsthm}
\usepackage{hyperref}

\newtheorem{theorem}{Theorem}
\newtheorem{lemma}{Lemma}
\newtheorem{remark}{Remark}

\newcommand{\R}{\mathbb R}
\newcommand{\N}{\mathbb N}
\newcommand{\Gb}{G_b^{0,\omega}}
\newcommand{\Cb}{C_b^{0,\omega}}
\newcommand{\MAP}{\mathrm{MAP}}

\title{A one-dimensional MAP subcase for geometric preduals of bounded
\texorpdfstring{\(C^{0,\omega}\)}{C0omega} spaces}
\author{Run \texttt{fa\_banach\_001}, lane 3}
\date{2026-06-21}

\begin{document}
\maketitle

\section*{Source and Question}

The source paper is Alexander Brudnyi, \emph{On Properties of Geometric
Preduals of \(C^{k,\omega}\) Spaces}, arXiv:1607.04824.  In Section 1.4,
after Theorem 1.3, the paper proves that \(\Gb(\R^n)\) has the metric
approximation property when \(\lim_{t\to\infty}\omega(t)=\infty\).  It then
observes that in the bounded-modulus case \(\Gb(\R^n)\) is merely isomorphic
to a space with the metric approximation property, with distortion depending
on \(\omega\), and states that it is not clear whether
\(G_b^{k,\omega}(\R^n)\) itself has the metric approximation property.

This packet proves a scoped positive result for the full-space case
\(k=0,n=1\), under the additional natural hypothesis that the modulus is
concave.

\section*{Proof Intuition}

For \(k=0\), the dual space is the bounded \(\omega\)-Lipschitz function
space: the unit ball consists of functions \(f\) with
\(\|f\|_\infty\le 1\) and
\(|f(x)-f(y)|\le \omega(|x-y|)\).  The key finite-rank approximants come from
randomly shifted grids.  If \(Q_{h,u}(x)=h\lfloor (x+u)/h\rfloor\) and
\(u\) is uniformly distributed on \([0,h]\), then for
\(t=mh+s\), \(0\le s<h\),
\[
 \mathbb E\,\omega(|Q_{h,u}(x+t)-Q_{h,u}(x)|)
 =
 \left(1-\frac{s}{h}\right)\omega(mh)
 + \frac{s}{h}\omega((m+1)h)
 \le \omega(t),
\]
because \(\omega\) is concave.  Thus averaging over shifted roundings is a
contraction for the \(\omega\)-Lipschitz seminorm.  Clipping the rounded
points to a large interval makes the range finite-dimensional while preserving
contractivity.  The mesh goes to zero and the clipping interval expands, so
the induced finite-rank contractions on the geometric predual converge to the
identity.

\begin{theorem}
Let \(\omega:[0,\infty)\to[0,\infty)\) be nondecreasing, concave, continuous,
and satisfy \(\omega(0)=0\).  Then the geometric predual
\(G_b^{0,\omega}(\R)\) of \(C_b^{0,\omega}(\R)\), in the sense of
arXiv:1607.04824, has the metric approximation property.
\end{theorem}

\begin{proof}
For \(k=0\), the norm on \(C_b^{0,\omega}(\R)\) is
\[
 \|f\|=\max\left\{\sup_{x\in\R}|f(x)|,\,
 \sup_{x\ne y}\frac{|f(x)-f(y)|}{\omega(|x-y|)}\right\}.
\]
The geometric predual \(\Gb(\R)\) is the closed linear span of the evaluation
functionals \(\delta_x\) inside \(\Cb(\R)^*\), and by Brudnyi's Theorem 1.2
we may identify \(\Cb(\R)\) with \(\Gb(\R)^*\).

Fix \(R>0\) and \(h>0\).  Let
\[
 Q_{h,u}(x):=h\left\lfloor\frac{x+u}{h}\right\rfloor,\qquad 0\le u<h,
\]
and let \(C_R(t)=\max\{-R,\min\{t,R\}\}\).  Define
\[
 (T_{R,h}f)(x)
   :=\frac1h\int_0^h f\bigl(C_R(Q_{h,u}(x))\bigr)\,du,
   \qquad f\in \Cb(\R).
\]
The values \(C_R(Q_{h,u}(x))\) belong to the finite set
\[
 A_{R,h}:=\{-R,R\}\cup \{jh:\ -R<jh<R,\ j\in\mathbb Z\}.
\]
Hence \(T_{R,h}\) has finite-dimensional range, since
\[
 (T_{R,h}f)(x)=\sum_{a\in A_{R,h}}p_a(x)f(a)
\]
for suitable scalar coefficient functions \(p_a\).

We next prove \(\|T_{R,h}\|\le 1\).  The supremum norm is immediate:
\(\|T_{R,h}f\|_\infty\le \|f\|_\infty\).  For the seminorm, take
\(x<y\) and put \(t=y-x\).  Write \(t=mh+s\), with \(m\in\mathbb Z_+\) and
\(0\le s<h\).  As \(u\) ranges uniformly over \([0,h]\), the difference
\(Q_{h,u}(y)-Q_{h,u}(x)\) is \(mh\) with probability \(1-s/h\) and
\((m+1)h\) with probability \(s/h\).  Since \(C_R\) is nonexpansive on
\(\R\) and \(\omega\) is nondecreasing,
\[
\begin{aligned}
 \frac1h\int_0^h
   \omega\!\left(\left|C_R(Q_{h,u}(y))-C_R(Q_{h,u}(x))\right|\right)du
 &\le
 \left(1-\frac{s}{h}\right)\omega(mh)
   +\frac{s}{h}\omega((m+1)h)\\
 &\le \omega(mh+s)=\omega(|x-y|),
\end{aligned}
\]
where the last inequality is concavity of \(\omega\).  Therefore, for
\(f\in\Cb(\R)\),
\[
\begin{aligned}
 |T_{R,h}f(y)-T_{R,h}f(x)|
 &\le
 \frac1h\int_0^h
 |f(C_R(Q_{h,u}(y)))-f(C_R(Q_{h,u}(x)))|\,du\\
 &\le |f|_{\omega}\,\omega(|x-y|).
\end{aligned}
\]
Thus \(T_{R,h}\) is a weak-star continuous finite-rank contraction on
\(\Cb(\R)=\Gb(\R)^*\).  Let \(H_{R,h}\) be its preadjoint on \(\Gb(\R)\);
equivalently,
\[
 H_{R,h}\delta_x
   =\frac1h\int_0^h \delta_{C_R(Q_{h,u}(x))}\,du,
\]
which is a convex combination of evaluations at points of \(A_{R,h}\).
Hence \(H_{R,h}\) is a finite-rank contraction.

It remains to check convergence to the identity.  Fix \(x\in\R\).  If
\(R>|x|+1\) and \(0<h<1\), then clipping does not affect
\(Q_{h,u}(x)\), and \(|Q_{h,u}(x)-x|\le h\).  Therefore
\[
 \|H_{R,h}\delta_x-\delta_x\|
 \le \frac1h\int_0^h \|\delta_{Q_{h,u}(x)}-\delta_x\|\,du
 \le \omega(h)\to 0.
\]
The linear span of the evaluations \(\delta_x\) is dense in \(\Gb(\R)\), and
the operators \(H_{R,h}\) are uniformly bounded by one.  Taking, for example,
\(R=N\) and \(h=1/N\), we obtain finite-rank contractions converging strongly
to the identity on \(\Gb(\R)\), and hence uniformly on compact subsets.  This
is precisely the metric approximation property.
\end{proof}

\section*{Verification Notes and Scope}

The proof uses only the \(k=0,n=1\) structure and concavity of the modulus.
It does not solve the higher-dimensional problem, the \(k\ge 1\) cases, or
the nonconcave moduli allowed by the source paper's weaker hypotheses.  For
unbounded concave moduli the source already implies the metric approximation
property; the new content is the bounded concave-modulus case, where the
source only gives an isomorphic MAP comparison and a bounded approximation
constant above one.

The shifted-grid contraction is insensitive to boundedness of \(\omega\).
Examples covered include \(\omega(t)=t^\alpha\) for \(0<\alpha\le 1\) and
bounded concave variants such as \(\omega(t)=\min\{t^\alpha,1\}\).

\section*{Novelty Check}

The run indexes had no exact previous result for arXiv:1607.04824.  Searches
for the exact arXiv id, for the notation \(G_b^{k,\omega}\), for
\(C_b^{k,\omega}\) and the approximation property, and for snowflake or
concave-modulus Lipschitz-free MAP variants did not locate a later direct
answer.  The local registry has a separate literature-implied packet for
arXiv:1501.07036 covering the usual-metric real-line subcase of a
Lipschitz-free MAP problem; that does not give this bounded
\(\omega(|x-y|)\)-metric result by itself.

\section*{References}

\begin{itemize}
\item A. Brudnyi, \emph{On Properties of Geometric Preduals of
\(C^{k,\omega}\) Spaces}, arXiv:1607.04824.
\item E. Pernecka and R. J. Smith, \emph{The metric approximation property
and Lipschitz-free spaces over subsets of \(\mathbb R^N\)}, arXiv:1501.07036.
\end{itemize}

\section*{Human Review Recommendation}

Recommended review focus: verify that the shifted-grid averaging operator is
indeed weak-star continuous with the stated preadjoint, and that the concavity
inequality is the only extra hypothesis beyond the \(k=0,n=1\) source
setting.

\end{document}

